% Predictions. Brief: briefs/predictions.md. Claims carried: the entries of ../claims.yaml with
% section predictions. Prose: PH-W8.3, in the regime frame. Every prediction is conditional on the
% screen's model and names its observable and what would refute it; no result is stated here.
\section{Predictions}\label{sec:predictions}

% (1) The initial experiment.
The screen turns into one experiment that tests both stages with the same foils. % CL-Q01
Aluminium foil is tested bare, carbon-coated, and carrying a calcium carbonate film, a calcium phosphate (hydroxyapatite) film, a baked aluminophosphate film or a film made by one of the in-situ routes, in two tests: (a) immersion in a model alkaline slurry at the pH and calcium doses of the processing grid, and (b) anodic potential holds at \NumClassLo\ and at \NumHoldLow\ to \NumHoldHigh\ in a fluorine-free electrolyte, such as lithium perchlorate in a carbonate solvent, the salt of the solid-salt sink. % CL-Q01
The observables are the calcium, phosphorus and aluminium released into the immersion liquid or the electrolyte, the anodic current against the bare and carbon-coated foils, the gas evolved during the hold, and the composition of the film before and after each test. % CL-Q01

% (2) The conditional predictions, each with its observable and what would refute it.
If the screen's model holds, then a calcite film stays undissolved in the model slurry inside the passing region of Section~\ref{sec:res-dissolution} and dissolves at the low-dose corner, and it is not oxidised at the \NumClassLo\ hold, below the measured onset of lithium carbonate at \NumTierTwoLow; it is thermodynamically excluded\nacreExcludedMark\ above about \NumCarbonateCrossover, so at the higher hold it is expected to lose calcium with release of carbon dioxide. % CL-Q05
Calcium released inside the passing region, or carbon dioxide at the low hold, would refute the dissolution or the low-voltage reading; a carbonate film unchanged at the higher hold would not refute the thermodynamic exclusion but would show kinetic stabilisation, a property the screen does not test. % CL-Q05
If the screen's model holds, then a hydroxyapatite film is not oxidised at the \NumClassLo\ hold and is thermodynamically excluded\nacreExcludedMark\ at the \NumHoldLow\ to \NumHoldHigh\ hold, above its crossover of \NumHapCrossover, where it is expected to lose calcium with oxygen evolution and leave pyrophosphate behind; calcium loss at the low hold would refute the onset ladder, and an unchanged film at the higher hold would again indicate kinetic stabilisation. % CL-Q02
If the screen's model holds, then a condensed calcium phosphate or aluminium phosphate film is not excluded at the \NumHoldLow\ to \NumHoldHigh\ hold: the ladder sets no bound below the hold for the calcium metaphosphate glass and none at all for the aluminium phosphates, so the hold measures what the screen leaves open. % CL-Q03
Oxidation of a calcium metaphosphate film well below its upper bound of \NumMetaGlassOnset\ would refute the ladder, because the upper bound holds for every sink; the bound of \ce{Ca2P2O7}, \NumPyroOnset, lies inside the hold range, so its oxidation there would refute nothing, and oxidation of an aluminium phosphate film would be a measurement where the screen has none. % CL-Q03

% (3) What would decide the nucleation condition.
Whether any of these films can be grown on a template from the slurry is outside the experiment, and the screen states it only as a model statement. % CL-O04
A measured or computed template--nucleus energy of hydroxyapatite to about \NumPrecGammaSnMj\ near zero would decide the nucleation condition, and a measured transfer free energy of \ce{Ca^2+} into the cell solvent to about \NumPrecTransferCaKj\ would allow an oxidation pass of a calcium phase at the low class; nothing in this paper supplies either. % CL-O04
